\documentclass[10pt,a4paper]{amsart}
\usepackage[margin=19mm]{geometry}
\usepackage{amsmath,amssymb,mathtools}
\usepackage[scr=rsfs,cal=euler]{mathalfa}
\usepackage{xfrac}
\usepackage[unicode,hidelinks,bookmarks=false]{hyperref}
\newcommand{\ie}{i.e.\ }
\newcommand{\cf}{cf.\ }
\newcommand{\resp}{respectively\ }
\newcommand{\Subsection}{\subsection}
\providecommand{\nobreakdash}{\nobreak\hbox{-}}
\setlength{\emergencystretch}{2em}
\multlinegap0pt
\usepackage{bm}
\usepackage{bbm}
\usepackage{dsfont}
\usepackage{xspace}
\usepackage{cancel}
\usepackage{mathtools}
\usepackage{amsthm}
\makeatletter
\@ifundefined{theorem}{\newtheorem{theorem}{Theorem}}{}
\@ifundefined{lemma}{\newtheorem{lemma}[theorem]{Lemma}}{}
\makeatother
\newcommand{\R}{{\mathbb R}}
\newcommand{\Z}{{\mathbb Z}}
\newcommand{\vol}{\operatorname{vol}}
\title[Weyl formula improvement for product of Zoll manifolds]{Weyl formula improvement for product of Zoll manifolds}
\author{Yanfei Wang}
\date{Source: arXiv:2506.02299v1 (CC BY 4.0)}
\begin{document}
\maketitle

\noindent\emph{Source and licence.} Weyl formula improvement for product of Zoll manifolds. Version arXiv:2506.02299v1 (CC BY 4.0), distributed under the Creative Commons Attribution 4.0 International licence, as recorded in the arXiv licence metadata for this exact version and shown on the abstract page https://arxiv.org/abs/2506.02299v1. Full rights notice: licenses/arxiv-2506.02299-CC-BY-4.0.txt.\par\smallskip

\noindent\textit{Editorial scope.} Counted unit: the Lemma whose source label is \texttt{ZollLatticeCount} in the article (source0.tex lines 832--841, source label \texttt{ZollLatticeCount}), delivered together with the article's own complete original proof of it (source0.tex lines 842--905, one \texttt{proof} environment of 64 lines). No mathematics is written, repaired or re-derived: statement and proof are the article's own text line for line. Required context retained uncounted: zec (lines 220--226); ZLC (lines 827--829). Every definition, display and label of those lines is retained unchanged. Omitted from this package and not redistributed: 1--219, 227--826, 830--831, 906--1164. Nothing inside a retained excerpt is omitted or altered: every retained body is delivered exactly as the article's lines are written, so no omission receipt of any kind accompanies this package. Citations: the retained excerpts carry no citation command at all, so no bibliography entry of the article is transcribed and no reference is redistributed. Reference adaptation: none was needed and none was made; every reference inside the delivered unit resolves inside it, so no unresolved reference or citation is produced. Printed slips kept literal: no printed word, symbol, hypothesis or number of the counted statement or of its proof was altered. Layout and class: the article's own class is \texttt{amsart} with packages including latexsym, amsmath, amssymb, bbold, xcolor; the wrapper is the lane's pinned \texttt{amsart} editorial wrapper at 10pt on A4, which restates the theorem-like environments the retained text needs and adds the article's own macro definitions that the retained text needs (\texttt{\textbackslash R}, \texttt{\textbackslash Z}, \texttt{\textbackslash vol}), with extra wrapper packages (bm, bbm, dsfont, xspace, cancel, mathtools). The delivered numbering is the unit's own. Assets: no retained span contains a figure environment, a graphics-inclusion command, a PDF-page inclusion, a table or a TikZ picture; no image file is copied into this package, the package declares no asset dependency and no image file is redistributed with it. Licence: version arXiv:2506.02299v1 of this article is distributed under the Creative Commons Attribution 4.0 International licence, as recorded in the arXiv licence metadata for this exact version and shown on the abstract page https://arxiv.org/abs/2506.02299v1. A full rights notice is delivered at licenses/arxiv-2506.02299-CC-BY-4.0.txt. Characters: every retained excerpt is pure ASCII, so the delivered engine, XeLaTeX with its default Latin Modern text font, reports no missing character.
\begin{lemma}\label{zec}
    Z is still our Zoll manifold defined as above, the number of eigenvalues in the k'th cluster as a polynomial of $k+\frac{\alpha}{4}$ is 
    \begin{equation} \label{zolclu}
    P_{Z^d}(t) = C \cdot t^{d-1} + O(t^{d-3}) \qquad \text{ for $|t|$ large.}
    \end{equation}
    where $t=k + \frac{\alpha}{4}$ and C is a constant depend on Z.
\end{lemma}

\begin{equation}\label{ZLC}
    \hat{N}(\lambda) : = \sum_{\substack{m \in \Z_{\geq 0}^k \times \Z^{n-k} \\ |m+y| \le \lambda }} \prod_{i = 1}^k P_{Z_i} (m_i +\frac{\alpha_i}{4})
\end{equation}

\begin{lemma}\label{ZollLatticeCount}
    Let $\hat{N}(\lambda)$ defined by \eqref{ZLC}. Then we have 
    \begin{equation}
        \hat{N}(\lambda) = \frac{|B_{|d|}|}{(2\pi)^{|d|}} \vol(M) \lambda^{|d|} + O(\lambda^{|d|-1-\frac{n-1}{n+1}}).
    \end{equation}
    and 
    \begin{equation}
        \hat{N}(\lambda) -\hat{N}(\lambda- c/\lambda) = O (\lambda^{|d|-2}).
    \end{equation}
\end{lemma}

\begin{proof}
    By lemma \ref{zec}, we know that 
    \begin{equation}
        P_{Z_i}(m_i+\frac{\alpha_i}{4}) = C_i (m_i+\frac{\alpha_i}{4})^{d_i-1} + O ((m_i+\frac{\alpha_i}{4})^{d_i-3}). \quad for \quad m_i \ge M
    \end{equation}
    To use this polynomial expansion, we need to define 
    \begin{equation}
        \hat{N}_1(\lambda) : = \sum_{\substack{m \in \Z_{\geq M}^k \times \Z^{n-k} \\ |m+y| \le \lambda }} \prod_{i = 1}^k P_{Z_i} (m_i +\frac{\alpha_i}{4})
    \end{equation}
    First we need to show that the difference between $\hat{N}_1(\lambda)$ and $\hat{N}(\lambda)$ won't hurt us much and is of order $O(\lambda^{d-1})$.
    \begin{align*}
        \hat{N}(\lambda)-\hat{N}_1(\lambda) &= \sum_{\substack{m \in \Z_{\geq 0}^k \times \Z^{n-k} \\ |m+y| \le \lambda }} \prod_{i = 1}^k P_{Z_i} (m_i +\frac{\alpha_i}{4}) - \sum_{\substack{m \in \Z_{\geq M}^k \times \Z^{n-k} \\ |m+y| \le \lambda }} \prod_{i = 1}^k P_{Z_i} (m_i +\frac{\alpha_i}{4}) \\
        & = \sum_{\substack{m \in \bigcup_{i=1}^k(\{m_i \le M\}\bigcap\Z_{\geq 0}^k \times \Z^{n-k}) \\ |m+y| \le \lambda }} \prod_{i = 1}^k P_{Z_i} (m_i +\frac{\alpha_i}{4})
    \end{align*}
    Since the first k manifolds are all Zoll manifolds with dimension greater or equal than 2, if we only count the first M lattice points in one of those directions, then we will lose at least two dimension in our Weyl counting function, which is the following
    \[
    \hat{N}(\lambda)-\hat{N}_1(\lambda) = O(\lambda^{|d|-2})
    \]
    Now we rewrite the $\hat{N}_1$ as
    \[
    \hat{N}_1(\lambda) = \sum_{m \in \Z_{\geq M}^k \times \Z^{n-k} + y} \chi_D(|m|/\lambda) \prod_{i=1}^k P_{Z_i}(m_i)
    \]
    And we define the main part of $\hat{N}_1$ as $\hat{N}_2$
    \[
    \hat{N}_2(\lambda) = \sum_{m \in \Z_{\geq M}^k \times \Z^{n-k} + y} \chi_D(|m|/\lambda) \prod_{i=1}^k C_im_i^{d_i-1}
    \]
    Then we can use Lemma \ref{zec} to estimate the difference of $\hat{N}_1$ and $\hat{N}_2$, which is also of order $O(\lambda^{|d|-2})$
    \begin{align*}
    \hat{N}_1(\lambda)-\hat{N}_2(\lambda) &= \sum_{m \in \Z_{\geq M}^k \times \Z^{n-k} + y} \chi_D(|m|/\lambda) \prod_{i=1}^k \left( C_i m_i^{d_i-1} + O(m_i^{d_i-3})\right)- \hat{N}_1(\lambda)\\
    &= \sum_{m \in \Z_{\geq M}^k \times \Z^{n-k} + y} \chi_D(|m|/\lambda) \prod_{i=1}^k O(|m|^{|d|-n-2})\\
    &= O(\lambda^{|d|-2})
    \end{align*}
    Where the last step is just due to the fact that the number of lattice points inside the $\lambda$-radius ball $B_{\lambda}$ is just $O(\lambda^n)$.

    Now we need to get rid of the restriction $m_i \geq M$ such that it's easy to compute and use the proposition below. So we define $\hat{N}_3(\lambda)$ as
    \[
    \hat{N}_3(\lambda) = \sum_{m \in  \Z^{n} + y \bigcap \R_{+}^k \times \R^{n-k}} \chi_D(|m|/\lambda) \prod_{i=1}^k C_im_i^{d_i-1}
    \]
    The difference between $\hat{N}_2(\lambda)$ and $\hat{N}_3(\lambda)$ is just
    \[
    \hat{N}_3(\lambda) - \hat{N}_2(\lambda) = \sum_{m \in  (\Z^{n} + y) \bigcap \bigcup_{j=1}^k H_j} \chi_D(|m|/\lambda) \prod_{i=1}^k C_im_i^{d_i-1}
    \]
    where $H_j = \{ x \in \R_{+}^k \times \R^{n-k} : x_j \in [0,y_j+M)\}$, for $j = 1, \cdots, k$. Actually the difference of $\hat{N}_3(\lambda) $ and $\hat{N}_2(\lambda) $ are very similar to the difference of $\hat{N}_1(\lambda) $ and $\hat{N}(\lambda) $, we could estimate them in the same way. We will elaborate the estimate of the first one because it is easier.
    \[
    \prod_{i=1}^k C_i m_i^{d_i-1} = O(|m|^{|d|-n-d_j+1}) \quad x \in H_j
    \]
    Hence 
    \[
    \sum_{m \in  (\Z^{n} + y) \bigcap H_j} \chi_D(|m|/\lambda) \prod_{i=1}^k C_im_i^{d_i-1} = O(\lambda^{|d|-d_j})
    \]
    just by naive estimate that the number of lattice point $\sum_{m \in  (\Z^{n} + y) \bigcap H_j} \chi_D(|m|/\lambda)$ is of $O(\lambda^{n-1})$. Since all $d_j \geq 2$, $O(\lambda^{|d|-d_j})$ is better than $O(\lambda^{|d|-2})$, so we have 
    \begin{equation}
        \hat{N}_3(\lambda) - \hat{N}_2(\lambda) = O(\lambda^{|d|-2})
    \end{equation}
    Combine all the above reduction, we have 
    \begin{equation}
        \hat{N}(\lambda) - \hat{N}_3(\lambda) = O(\lambda^{|d|-2})
    \end{equation}
    If we could prove that $\hat{N}_3(\lambda)$ satisfy
    \begin{equation}
        \hat{N}_3(\lambda) = \frac{|B_{|d|}|}{(2\pi)^{|d|}} \vol(M) \lambda^{|d|} + O(\lambda^{|d|-1-\frac{n-1}{n+1}}).
    \end{equation}
    We can finish the proof. Now what we need is just to prove the following proposition.
\end{proof}

\end{document}
